\section{Resumen final: qué memorizar para la etapa digital}

\begin{memorizar}
\begin{multicols}{2}
\textbf{Definiciones exactas:}
\begin{itemize}[leftmargin=1.2em, itemsep=0pt]
  \item MPS = cuántos productos; MRP = qué materiales y cuándo.
  \item L4L minimiza inventario, NO setups.
  \item Holgura = LS--ES = LF--EF; RC tiene H=0.
  \item Kingman: CT $= V\cdot U\cdot T$.
  \item Little: $L=\lambda W$; WIP$=$TH$\cdot$CT.
  \item OEE $=$ Disp $\times$ Rend $\times$ Calidad.
\end{itemize}

\textbf{Diferencias clave:}
\begin{itemize}[leftmargin=1.2em, itemsep=0pt]
  \item Chase (ajusta capacidad) vs.\ Nivel (inventario/faltante).
  \item Explotar (sin \$) vs.\ Elevar (con \$).
  \item Cross-dock (sin stock) vs.\ Bodega (almacena).
  \item Postponement reduce inventario; Risk pooling reduce SS.
\end{itemize}

\columnbreak

\textbf{La Meta 16--26 --- 5 ideas:}
\begin{enumerate}[leftmargin=1.2em, itemsep=0pt]
  \item 5 pasos: Identificar, Explotar, Subordinar, Elevar, Repetir.
  \item DBR: Tambor=CB, Buffer antes del CB, Cuerda ata liberación.
  \item Inercia es el enemigo (volver al paso 1).
  \item Gestionar = método socrático (preguntar).
  \item Lotes pequeños $\Rightarrow$ mejor flujo.
\end{enumerate}

\textbf{Fórmulas a recordar:}
\begin{itemize}[leftmargin=1.2em, itemsep=0pt]
  \item $t_e = \tfrac{a+4m+b}{6}$,\ $\sigma=\tfrac{b-a}{6}$
  \item $\rho=\lambda/\mu$,\ $L_q=\tfrac{\rho^2}{1-\rho}$
  \item Slope $=\tfrac{CC-NC}{NT-CT}$
  \item Silver-Meal: parar si $CT(k)\!>\!CT(k\!-\!1)$
\end{itemize}
\end{multicols}
\end{memorizar}

\begin{paradesarrollo}
\begin{itemize}[leftmargin=1.2em, itemsep=1pt]
  \item \textbf{Lleva este formulario impreso} (Sección 9).
  \item Ejercicios casi seguros: \textbf{modelo LP de Planif.\ Agregada}, \textbf{PERT con probabilidad/bono}, \textbf{MRP (BOM + L4L + Silver-Meal)}, \textbf{Colas (M/M/1, G/G/1)}.
  \item En modelamiento: define SIEMPRE conjuntos, parámetros, variables, F.O.\ y restricciones por separado. Las binarias para condiciones lógicas dan puntos fáciles.
  \item En MRP: cuidado con el desfase del lead time y con que la orden del padre genera el GR del hijo.
  \item En PERT: identifica bien la ruta crítica antes de calcular varianzas; usa la tabla normal para bono/penalización.
  \item En Colas: verifica $\rho<1$; recuerda que Kingman es VUT.
  \item \textbf{Planificación Agregada, MRP y Proyectos} son los temas más probables en desarrollo (confirmado por I2 2025).
\end{itemize}
\end{paradesarrollo}

